\documentclass[10pt,a4paper]{amsart}
\usepackage[margin=19mm]{geometry}
\usepackage{amsmath,amssymb,mathtools}
\usepackage[scr=rsfs,cal=euler]{mathalfa}
\usepackage{xfrac}
\usepackage[unicode,hidelinks,bookmarks=false]{hyperref}
\newcommand{\ie}{i.e.\ }
\newcommand{\cf}{cf.\ }
\newcommand{\resp}{respectively\ }
\newcommand{\Subsection}{\subsection}
\providecommand{\nobreakdash}{\nobreak\hbox{-}}
\setlength{\emergencystretch}{2em}
\multlinegap0pt
\usepackage{amssymb}
\usepackage{mathtools}
\usepackage{xcolor}
\usepackage[shortlabels]{enumitem}
\usepackage{algorithm}\usepackage{algpseudocode}
\usepackage{tikz}
\newtheorem{lemma}{Lemma}
\newtheorem{theorem}{Theorem}
\title{Sequential KV Cache Compression\\via Probabilistic Language Tries:\\ Beyond the Per-Vector Shannon Limit}
\author{Gregory Magarshak}
\date{Source: arXiv:2604.15356v1 (CC BY 4.0)}
\begin{document}
\maketitle

\noindent\emph{Source and licence.} Sequential KV Cache Compression\\via Probabilistic Language Tries:\\ Beyond the Per-Vector Shannon Limit. Version arXiv:2604.15356v1 (CC BY 4.0), distributed by arXiv under the Creative Commons Attribution 4.0 International licence, http://creativecommons.org/licenses/by/4.0/, as recorded in the arXiv licence metadata for this exact version and shown on the abstract page https://arxiv.org/abs/2604.15356v1. Full licence notice: \texttt{licenses/arxiv-2604.15356-CC-BY-4.0.txt}.\par\smallskip

\noindent\textit{Editorial scope.} Counted unit: the article's theorem thm:seq\_bound (source0.tex lines 343--357). The article's complete original proof of it is delivered with it (source0.tex lines 359--413, one \texttt{proof} environment of 55 lines). No mathematics is written, repaired or re-derived: the statement and the proof are the article's own lines, in the article's own notation, and no retained line is substituted or re-flowed.  Retained uncounted as the source-local context the proof needs: lines 313--321.  Nothing was omitted from any retained span, so the package carries no omission record.  No retained line contains a citation, so the wrapper carries no bibliography and no bibliography entry.  No retained line contains a figure environment, an image-inclusion command or drawing code, so no image is delivered and the package declares no asset dependency.  Editorial formatting only, in addition: an AMS \texttt{amsart} preamble that restates the article's own declarations and macros used by the retained lines, namely the environments \texttt{lemma}, \texttt{theorem} on separate counters, together with the standard packages those lines need.  The retained environments print in this document as Lemma 1, Theorem 1 in order of appearance, whereas the article prints the same items with the numbering of its own full document.  The complete native source of the article as distributed in the arXiv e-print for this exact version is bundled unchanged as \texttt{sources/198600/source0.tex}.
\begin{lemma}[KV determinism]
\label{lem:determinism}
For a fixed model $\mathcal{M}$ with deterministic forward pass, $\mathrm{KV}_i$
is a deterministic function of the token prefix $(t_1, \ldots, t_i)$:
\[
  \mathrm{KV}_i = F_{\mathcal{M}}(t_1, \ldots, t_i)
\]
for some deterministic function $F_{\mathcal{M}}: V^i \to \mathbb{R}^{2LH_{\mathrm{head}}d}$.
\end{lemma}

\begin{theorem}[Sequential entropy bound]
\label{thm:seq_bound}
Let $\mathbf{t} = (t_1, t_2, \ldots)$ be a random token sequence drawn from
$P_{\mathcal{M}}$.  For any position $i \geq 2$ (the bound holds trivially at $i=1$ since
$H(\mathrm{KV}_1) \leq H(t_1)$ by Lemma~\ref{lem:determinism} and data processing):
\[
  H\bigl(\mathrm{KV}_i \mid \mathrm{KV}_{\leq i-1}\bigr)
  \;\leq\;
  H\bigl(t_i \mid t_1, \ldots, t_{i-1}\bigr)
  \;=\;
  H\bigl(t_i \mid \mathrm{KV}_{\leq i-1}\bigr).
\]
The conditional entropy of the $i$-th KV vector, given all prior cache entries,
is bounded above by the model's per-token surprisal at position $i$.
\end{theorem}

\begin{proof}
The proof proceeds in two steps: first establishing that $\sigma(\mathrm{KV}_{\leq i-1}) = \sigma(t_{\leq i-1})$, then deriving the bound.

\textbf{Step 1: Injectivity and sigma-algebra equivalence.}

We establish that $t_j$ is a measurable function of $\mathrm{KV}_{\leq j}$ for
each $j$, by showing that the map $t_j \mapsto \mathrm{KV}_j \mid t_{<j}$ is
injective.

Fix any context $t_{<j}$ and suppose $t_j \neq t_j'$.  At layer $\ell = 1$,
the key vector is $\mathbf{k}^{(1)}_j = W_K^{(1)} E(t_j)$.  Since the
embedding matrix $E: V \to \mathbb{R}^{d_{\rm model}}$ has full column rank
in any trained transformer (distinct tokens receive distinct embeddings; the
set of weight matrices for which any two token embeddings coincide has measure
zero in parameter space), we have $E(t_j) \neq E(t_j')$, and since
$W_K^{(1)} \in \mathbb{R}^{d_{\rm head} \times d_{\rm model}}$ satisfies the generic condition that no pairwise embedding difference $E(t) - E(t')$ lies in $\ker W_K^{(1)}$ (this holds with probability~1 over random weight matrices since the null space has dimension $d_{\rm model} - d_{\rm head}$ and there are finitely many token pairs), it follows that
$\mathbf{k}^{(1)}_j(t_j) \neq \mathbf{k}^{(1)}_j(t_j')$.  In particular,
$\mathrm{KV}_j^{(1)}(t_j) \neq \mathrm{KV}_j^{(1)}(t_j')$.

Since $\mathrm{KV}_j$ includes the layer-1 component $\mathrm{KV}_j^{(1)}$
as a subvector, the full vector $\mathrm{KV}_j$ distinguishes $t_j$ from $t_j'$.
Thus $t_j \mapsto \mathrm{KV}_j \mid t_{<j}$ is injective for each $j$.
(No assumption about higher-layer projections is needed: injectivity is
witnessed by the layer-1 keys alone.)

By Lemma~\ref{lem:determinism}, $\mathrm{KV}_j$ is measurable w.r.t.\ $\sigma(t_{\leq j})$
for each $j$.  By the injectivity just established, $t_j$ is measurable w.r.t.\
$\sigma(\mathrm{KV}_j, t_{<j})$.  Applying this inductively:
$t_1$ is determined by $\mathrm{KV}_1^{(1)}$ (hence by $\mathrm{KV}_1$);
$t_2$ is determined by $(\mathrm{KV}_2, t_1)$, hence by $(\mathrm{KV}_1, \mathrm{KV}_2)$;
and so on.
Therefore $t_{<i}$ are jointly measurable w.r.t.\
$\sigma(\mathrm{KV}_1,\ldots,\mathrm{KV}_{i-1})$, giving:
\begin{equation}
  \sigma(\mathrm{KV}_{\leq i-1}) = \sigma(t_{<i}).
  \label{eq:sigma_equiv}
\end{equation}

\textbf{Step 2: The bound.}

Since $\mathrm{KV}_i = F_{\mathcal{M}}t_{\leq i}$ is a deterministic function
of $(t_{\leq i-1}, t_i)$, the data-processing inequality gives:
\[
  H(\mathrm{KV}_i \mid t_{\leq i-1}) \leq H(t_i \mid t_{\leq i-1}).
\]
Applying~\eqref{eq:sigma_equiv} to both sides (conditioning on $\mathrm{KV}_{\leq i-1}$
is equivalent to conditioning on $t_{\leq i-1}$):
\[
  H(\mathrm{KV}_i \mid \mathrm{KV}_{\leq i-1})
  = H(\mathrm{KV}_i \mid t_{\leq i-1})
  \leq H(t_i \mid t_{\leq i-1})
  = H(t_i \mid \mathrm{KV}_{\leq i-1}),
\]
which is the stated inequality and equality simultaneously.
\end{proof}

\end{document}
